Prior work established that reviewer profile depth is the dominant predictor of AI evaluation quality for puzzle hunt puzzles, with a non-monotonic gradient in which design philosophy captures most of the value of a full operational profile. Whether the same depth gradient governs AI \emph{authoring} --- and whether reviewing expertise transfers to creative generation --- is an open question, because reviewing applies a lens to finished material while authoring must generate material that does not yet exist.

We conducted a 7-condition authoring ablation (A0--A6) mirroring a prior reviewer ablation, extended with a 3-designer $\times$ 2-mode comparison in which Snyder, Katz, and Young each authored puzzles under both a full authoring profile (A5) and their own C8-tier reviewing profile repurposed as authoring guidance (A6). All 13 conditions used a fixed knowledge-extraction task (Age of Empires~II, answer \textsc{tower}, five-clue first-letter acrostic) evaluated by a same-session C8 expert panel on a 7-dimension weighted rubric (/45).

Three findings emerge. First, the authoring quality gradient mirrors the reviewer ablation: improvement is real but non-monotonic, with domain world data alone adding negligible value over bare task specification (A0~=~21.7, A2~=~28.3), and the largest single jump occurring at principle introduction (A2 to A3, $+2.7$ points). Second, transfer quality depends on lens type: operational lenses (Katz, $+1.6$ points) invert cleanly into authoring-time constraints; construction lenses (Snyder, $-0.7$) are already generatively oriented and require no transposition; perceptual lenses (Young, $+3.7$) are format-dependent --- when the output format cannot supply the required visual material, the reviewer-as-author exercise forces a productive structural substitution, yielding the dataset high of 38.0/45. Third, a format ceiling constrains all A-series puzzles to Riven Standard $\leq 3$: profile depth raises the quality floor within a fixed mechanism, but mechanism choice determines the ceiling.

The paper's primary contribution is a lens-type transfer model classifying reviewer profiles by generativity --- construction, operational, perceptual --- that predicts whether repurposing a reviewing profile for authoring will improve, preserve, or require reformulation. Empirical confirmation of the symmetric gradient establishes that the structural properties governing profile depth are not evaluation-specific. Practical guidance is given for constructing authoring profiles from reviewing profiles by lens type.
